\subsection{The cuspidal fibre and finite monodromy}
\label{sec:cusp}

\begin{proposition}[Rooting a fixed boundary cannot repair a lattice kernel;
\status{Proved}]\label{P1e-roots}
Suppose the normal-meridian map $\rho:\Z^r\to\pi _1(U)$ of a
normal-crossing pair has nonzero kernel. Taking arbitrary finite roots
of those same boundary components does not make this map injective
on a root-stack atlas. The assertion holds both discretely and
profinally.
\end{proposition}
\begin{proof}
For root indices $n_1,\ldots,n_r$, an atlas has coarse coordinates
$x_i=z_i^{n_i}$. Its meridian map is
$\rho\circ\operatorname{diag}(n_1,\ldots,n_r)$. The subgroup
$\bigoplus_i n_i\Z$ has finite index in $\Z^r$ and meets every
nonzero subgroup nontrivially: multiply any nonzero kernel vector
by a common multiple of the $n_i$. Thus the new map still has a
nonzero integral kernel. The same integral vector is nonzero in
$\widehat\Z^r$, and gives the profinite assertion. More generally,
if a purely profinite kernel vector $v$ is nonzero, multiplication
by a common root index $N$ puts it in the root lattice and cannot
kill it: $\widehat\Z^r$ has no nonzero integral torsion. This is a statement
about rooting an unchanged resolved boundary, not about replacing
the degeneration by a different stable model.
\end{proof}

We distinguish the affine surface $y^2=x^3+t$ from its projective
Weierstrass completion in the fibre direction. Both interpretations
of the cuspidal-fibre example can be computed completely.

\begin{proposition}[The resolved affine cusp; \status{Proved}]
\label{P1e-cusp}
Let $f(x,y)=y^2-x^3$, let $U=\C^2\setminus f^{-1}(0)$, and resolve
the reduced cusp divisor minimally to normal crossings. The resulting
pair has eight connected strata: the open $U$, four curve strata
$C^\circ,E_1^\circ,E_2^\circ,E_3^\circ$, and three nodes
$\nu_i=C_i\cap E_3$, where $C_1=E_1,C_2=E_2,C_C=C$.
They are
\[
 C^\circ=\C^*,\qquad E_1^\circ=E_2^\circ=\mathbb A^1,\qquad
 E_3^\circ=\PP^1\setminus\{0,1,\infty\}.
\]
Every stratum is aspherical. Put
\[
 G=\pi _1(U)=\langle a,b\mid a^2=b^3\rangle,\qquad
 h=a^2=b^3,\qquad m=ab^{-1}=h(ba)^{-1}.
\]
The meridians of $C,E_1,E_2,E_3$ are respectively $m,b,a,h$.
All four have infinite order. The full node maps have kernels
\[
 \ker\bigl(\Z\langle b\rangle\oplus\Z\langle h\rangle\to G\bigr)
       =\Z(-3,1),\qquad
 \ker\bigl(\Z\langle a\rangle\oplus\Z\langle h\rangle\to G\bigr)
       =\Z(-2,1),
\]
whereas $\Z\langle m\rangle\oplus\Z\langle h\rangle\to G$ is injective.
Thus the ordinary-exit-category criterion fails at exactly the two
exceptional-tail nodes.
\end{proposition}
\begin{proof}
The three relevant blow-up charts give, successively,
\[
 f=x^2(v^2-x),\qquad
 f=u^2v^3(v-u),\qquad
 f=u^6w^3(w-1).
\]
The last exceptional component meets the strict cusp and the two
earlier exceptionals at three distinct points. The multiplicities
of the total fibre on $(C,E_1,E_2,E_3)$ are $(1,2,3,6)$, and the
exceptional self-intersections are $(-3,-2,-1)$. The strict affine
cusp is its normalization with the origin removed, hence $\C^*$;
each early exceptional has just one puncture.

The weighted radial action $(x,y)\mapsto(r^2x,r^3y)$ identifies the
complement's homotopy type with the complement of the $(2,3)$ torus
knot. Splitting the knot exterior into its two solid-torus pieces
and applying van Kampen identifies their core powers $a^2$ and
$b^3$, giving the displayed presentation. Equivalently, the base
change $t=s^6$, $x=s^2X$, $y=s^3Y$ trivializes the fibration as
\[
 F\times\C^*,\qquad F=\{Y^2=X^3+1\}.
\]
The smooth projective completion of $F$ is elliptic and has one
point at infinity. Its affine fundamental group is free of rank
two, and the fibration over $\C^*$ proves that $U$ is aspherical.
The degree map induced by $f$ is
\[
 \varphi:G\to\Z,\qquad
 \varphi(a)=3,\quad\varphi(b)=2,\quad\varphi(h)=6,\quad\varphi(m)=1.
\]
The meridian identification can also be read without any ambiguity
from the resolution plumbing. A rational tail of square $-d$ imposes
$\mu_{\rm adjacent}=\mu_{\rm tail}^{d}$: this is the Euler-number
relation in the circle bundle of its normal line. Thus the two tails
impose $b^3=h$ and $a^2=h$. The central component has square $-1$,
so its three peripheral meridians satisfy $bam=h$, with a compatible
cyclic order and normal framing. Hence $m=h(ba)^{-1}=ab^{-1}$.
This identifies the chosen presentation and all four meridians
simultaneously. Their displayed degrees prove that all are infinite.
At the two tail nodes the images are the cyclic groups generated by
$b$ and $a$, with the exact kernels stated above. At the cusp node,
pass to $G/\langle h\rangle=C_2*C_3$. The image of $m=ab^{-1}$ is a
cyclically reduced word of positive alternating length, so no
nonzero power is trivial. A relation $m^rh^s=1$ consequently has
$r=0$, and then $s=0$ by $\varphi(h)=6$. The curve and point
strata have the claimed asphericity directly.
\end{proof}

\paragraph{All objects, mapping spaces, actions, and compositions.}
Here is a concrete group-and-link model for the entire exit category,
not merely for its nontruncated homotopy groups. There is one object
for each of the eight connected strata. The endomorphism groups
in the homotopy category are
\[
 G,\quad \Z\ (C^\circ),\quad 1\ (E_1^\circ),\quad
 1\ (E_2^\circ),\quad F_2\ (E_3^\circ),\quad 1,1,1.
\]
These are the components of the endomorphism mapping spaces; those
spaces themselves are homotopy-discrete. For each curve stratum
$S$ with meridian $\mu_S$, its exit space to $U$ is the discrete
coset set $G/\langle\mu_S\rangle$. The open group acts on the left.
For $C^\circ$ choose the normal framing in which the generator
of its stratum group lifts to $h$; this specifies its right action
on $G/\langle m\rangle$. For $E_3^\circ$ write
\[
 F_2=\langle p_1,p_2\rangle,\qquad p_3=(p_1p_2)^{-1}.
\]
The right action on $G/\langle h\rangle$ is induced by
$p_1\mapsto b,\ p_2\mapsto a$. The tail source groups are trivial.

The remaining nonempty off-diagonal spaces are
\[
\begin{array}{c|c}
\text{exit space}&\text{homotopy type}\\ \hline
\Map(\nu_1,U)&\displaystyle\coprod_{G/\langle b\rangle}S^1\\
\Map(\nu_2,U)&\displaystyle\coprod_{G/\langle a\rangle}S^1\\
\Map(\nu_C,U)&G/\langle m,h\rangle\\
\Map(\nu_1,E_1^\circ),\ \Map(\nu_2,E_2^\circ)&S^1\\
\Map(\nu_C,C^\circ)&*\\
\Map(\nu_i,E_3^\circ)&F_2/\langle p_i\rangle\quad(i=1,2,3).
\end{array}
\]
All other off-diagonal spaces are empty. To specify every face map
and composition, use the following link-group diagram and take
nerves and homotopy fibres, as in the link reconstruction in Part~1.
The punctured normal-link groups of the curves and their maps to $G$
are
\[
\begin{array}{c|c|c}
S&\pi _1(P_S)&\pi _1(P_S)\to G\\ \hline
E_1^\circ&\Z\langle b\rangle&b\mapsto b\\
E_2^\circ&\Z\langle a\rangle&a\mapsto a\\
C^\circ&\Z\langle m,h\rangle&(m,h)\mapsto(m,h)\\
E_3^\circ&\Z\langle h\rangle\times F_2&
             (h,p_1,p_2)\mapsto(h,b,a).
\end{array}
\]
At a node take its full rank-two group, with ordered generators
$(\mu_{\rm leaf},h)$. Its map to the leaf link is respectively
\[
 (b,h)\mapsto(b,b^3),\qquad
 (a,h)\mapsto(a,a^2),\qquad
 (m,h)\mapsto(m,h).
\]
Its map to the central link sends these generators respectively to
\[
 (p_1,h),\qquad(p_2,h),\qquad(hp_3,h).
\]
Projection from a curve link to its stratum kills its normal
generator. These formulas commute with the displayed maps to $G$;
in particular $hp_3\mapsto h(ba)^{-1}=m$. The nerves, their
homotopy fibres, and the induced maps for the two routes from a
node to $U$ specify a simplicial model of the category, including
composition homotopies. The tail-circle generator maps to
$(-3,1)$ or $(-2,1)$ in the corresponding node-to-open space.
Thus neither these circles nor their composition data may be
discarded. The displayed mapping spaces also follow directly from
the homotopy-fibre formula: for a homomorphism of discrete groups,
each component of its classifying-space fibre is the classifying
space of its kernel.

\begin{proposition}[The projective Weierstrass interpretation;
\status{Proved}]\label{P1e-cusp-proper}
For the projective family $y^2z=x^3+tz^3$, the same resolution has
$C^\circ=\mathbb A^1$ instead of $\C^*$; all three leaves are affine lines.
Its open fundamental group is
\[
 G_{\rm pr}=\Z^2\rtimes_A\Z,\qquad
 A=\begin{pmatrix}0&-1\\1&1\end{pmatrix},\quad A^6=1.
\]
It is the quotient of $G$ by the normal closure of
$\lambda=hm^{-6}$. The three node kernels, in the leaf-central
ordering, are now $\Z(-3,1)$, $\Z(-2,1)$, and $\Z(-6,1)$.
Every stratum is still aspherical, but all three full node lattices
fail injectivity. In the previous category replace $G$ by
$G_{\rm pr}$, replace the $C^\circ$ endomorphism group by $1$,
replace $\Map(\nu_C,C^\circ)=*$ by $S^1$, and replace
$\Map(\nu_C,U)$ by $\coprod_{G_{\rm pr}/\langle m\rangle}S^1$.
The third leaf-link homomorphism becomes
$(m,h)\mapsto(m,m^6)$. These replacements, with the same central
link diagram, give the entire projective exit category.
\end{proposition}
\begin{proof}
Filling the unique puncture of $F$ kills its peripheral commutator,
which is the knot longitude $\lambda=hm^{-6}$. Indeed the framing
generator $h$ goes once along the knot and has linking number six,
while $m$ is the positive meridian of linking number one; subtracting
six meridians gives the primitive zero-linking longitude. The remaining
fibre is a torus. The section at infinity splits the fundamental
group sequence, and its order-six monodromy acts on a suitable
homology basis by $A$. For an explicit compatible realization in
the semidirect product, with product
$(v,n)(w,k)=(v+A^nw,n+k)$, take
\[
 a=(e_2,3),\qquad b=(e_1,2),\qquad
 m=(0,1),\qquad h=(0,6).
\]
They satisfy $a^2=b^3=m^6=h$ and generate the group, since $a,b$
together with $m$ generate the two fibre basis vectors. The degree
map is still the second coordinate, so the three leaf meridians
are infinite, and the stated kernels follow from the three exact
power relations. The local resolution is unchanged. The proper
strict transform has square $-6$; the full fibre intersection
matrix annihilates the multiplicity vector $(1,2,3,6)$.
The torus fibration over the punctured disc is aspherical, and
the remaining strata were already aspherical. The homotopy-fibre
calculation of the preceding paragraph gives all the replacements.
\end{proof}

\begin{proposition}[The good orbifold replacement; \status{Proved}]
\label{P1e-cusp-stack}
Let $\zeta$ be a primitive sixth root of unity and let
\[
 \mathcal X=[F\times\Delta_s/\mu _6],\qquad
 \zeta\cdot(X,Y,s)=(\zeta^{-2}X,\zeta^{-3}Y,\zeta s).
\]
Its boundary is $\mathcal F=[F/\mu _6]$ at $s=0$, and its open
complement is the original affine family over $\Delta_t^*$,
where $t=s^6$, $x=s^2X$, and $y=s^3Y$.
The pair is good. Its boundary is the orbifold affine line with
cone orders $2,3$, so
\[
 \pi _1(\mathcal F)=C_2*C_3=G/\langle h\rangle .
\]
Its normal meridian is $h$. The exit category is ordinary up to
equivalence, with two objects, automorphism groups
$C_2*C_3$ and $G$, and crossing set $G/\langle h\rangle$.
The right action is that of the boundary quotient and the left
action is that of $G$.

Replacing $F$ by its projective elliptic completion gives the
corresponding good projective orbifold model. Its boundary group
is $\Z^2\rtimes_A\mu _6=G_{\rm pr}/\langle h\rangle$; its coarse
boundary is $\PP^1$ with cone orders $2,3,6$. The same two-object
description holds with these groups. This replacement is not
obtained by merely rooting the fixed three-blow-up NC divisor.
\end{proposition}
\begin{proof}
The action is free for $s\ne0$. The displayed invariant formulas
identify its quotient there with the given family. The affine
quotient coordinate on the boundary is $w=Y^2=X^3+1$.
At $w=0$ the stabilizer has order two; at $w=1$ it has order
three. The projective completion adds the fixed point at infinity,
of order six. The finite atlases $F$ and $\overline F$ are
aspherical, so their orbifold classifying spaces are aspherical
as well: the Borel fibrations over $B\mu _6$ have no higher
homotopy groups. A circle in the atlas normal coordinate $s$
is the central fibre of this punctured normal circle bundle.
In the affine group the centre is $\langle h\rangle$, since
its quotient $C_2*C_3$ is centreless. In the proper group the
centre is $\langle(0,6)\rangle$, since $A$ has no nonzero fixed
vector. The circle has degree six in the $t$-base, so it is
exactly $h$ and has infinite order. Its full profinite completion
also injects, since the composite with the completed degree map is
multiplication by six on $\widehat\Z$.

The punctured normal-link group is $G$ (respectively $G_{\rm pr}$);
its map to the boundary group is quotient by $\langle h\rangle$,
and its map to the open group is the identity. Taking the
homotopy fibre of the resulting graph homomorphism gives the
discrete crossing set asserted above, with precisely the stated
actions and group multiplication compositions. The criterion
from Part~1 applies.

The affine family is of stable pointed type $(1,1)$ after adding
its missing point. In the projective interpretation that point
can remain a marking rather than a removed boundary component.
We do not call unpointed genus one a stable Deligne--Mumford
curve. The construction is a base change, replacement by the
smooth elliptic model, and quotient descent. Proposition
\ref{P1e-roots} proves separately why finite roots of the old
resolved boundary cannot have the same effect.
\end{proof}

\begin{proposition}[What finite monodromy does and does not imply;
\status{Proved}]\label{P1e-finite}
Finite nontrivial monodromy does not by itself force failure of
meridian injectivity. For example, let a finite cyclic group act
freely on an aspherical smooth curve $F$ and act on a disc by
rotation. Then $(F\times\Delta)/\mu_n$ is a smooth manifold,
its boundary is the aspherical curve $F/\mu_n$, and the normal
meridian has degree $n$ in the punctured-disc base, so the pair
is good. The cusp obstruction above is the exact dependence
among several resolved meridians. Stable orbifold replacement
retains the finite monodromy in the boundary stabilizers while
retaining the infinite normal meridian in the open fundamental
group. All assertions here are over $\C$ or in characteristic
zero; no claim about wild reduction at $2$ or $3$ is made.
\end{proposition}
\begin{proof}
Freeness on $F$ makes the diagonal quotient a manifold even at
$s=0$. The finite covers $F\to F/\mu_n$ and
$F\times\Delta^*\to (F\times\Delta^*)/\mu_n$
identify their higher homotopy groups with those of aspherical
spaces, so both strata are aspherical. In a local normal coordinate
the invariant base parameter is $t=s^n$; the normal circle has degree
$n$ under the map of the open complement to $\Delta_t^*$.
Its cyclic fundamental group therefore injects. The remaining
statements follow from the preceding three cusp calculations and
the fixed-boundary root obstruction.
\end{proof}
